\PostTableHeadingSpace
\section{Discussion}

The validation results characterize the implemented architecture at the systems-engineering layer. In the bounded public-data slice, source grounding, entity routing, answer structure, trace generation, and output hygiene are represented as explicit artifacts and exercised through repeatable checks. The discussion below explains what this means for productization and reproducibility.

\subsection{Engineering Interpretation}

The reconstruction changes the unit of control from prompt instructions to versioned engineering artifacts. The motivating PoC surfaced a productization pattern in which useful behavior appeared before the system had explicit source, trace, and validation contracts. In the reconstructed system, source manifests, company-scoped claims, code-owned routing, answer contracts, and validation traces govern the answer process. This approach is complementary to prompt engineering: prompts still guide language behavior, while the surrounding architecture assigns ownership for evidence boundaries, entity scope, output structure, leakage prevention, and inspection.

The proposed pattern also differs from adjacent engineering approaches in emphasis. DSPy-style program optimization focuses on compiling and optimizing language-model pipelines \citep{khattab2024dspy}, while MLOps addresses broader lifecycle deployment concerns \citep{kreuzberger2023mlops}. The architecture described here is narrower: it defines how a generated enterprise answer remains connected to source manifests, source-backed claims, routing metadata, answer contracts, and traces. This distinction matters because enterprise readers need a useful answer, while developers and researchers need a separate trace that explains how the answer was assembled.

The results have three implications for the implemented system. First, the fixed validation-scenario runs show that expected source, routing, trace, and answer-structure contracts were preserved across the reference slice. Second, the fault-injection results show that the validators detected deliberately broken source, routing, trace, answer, leakage, link, and latency conditions. Third, the latency result shows that orchestration and caching improved runtime behavior and preserved the answer contract. The latency result is an operational signal; the broader contribution is the architecture-level separation of sources, claims, routing, answer contracts, traces, and validation gates.

\subsection{Productization Implications}
\label{sec:reproducibility}
\enlargethispage{3\baselineskip}

In practitioner discussions, organizations often ask how to locate their current AI-agent adoption state before full deployment. The harness becomes useful when product behavior must move from plausible interaction to governed operation: source scope, entity routing, claim eligibility, answer contracts, and traces become explicit organizational capabilities. \Cref{tab:maturity-ladder} translates that productization need into a descriptive adoption framework: Level 1 is prompt-centered, while Level 5 represents a traceable system with governed evidence and validation artifacts. The framework supports planning; scoring a particular organization remains a separate assessment.

\begin{table}[H]
    \centering
    \caption{Descriptive adoption framework for enterprise LLM-agent systems.}
    \label{tab:maturity-ladder}
    \footnotesize
    \renewcommand{\arraystretch}{1.12}
    \begin{tabular}{C{0.12\linewidth} L{0.32\linewidth} L{0.44\linewidth}}
        \toprule
        \multicolumn{1}{C{0.12\linewidth}}{Level} & \multicolumn{1}{C{0.32\linewidth}}{System state} & \multicolumn{1}{C{0.44\linewidth}}{Observable capability} \\
        \midrule
        1 & Prompt-centered chatbot & Product behavior is mainly controlled by instructions; source scope and validation remain informal. \\
        2 & RAG-based assistant & Retrieved documents support answers; evidence boundaries and traces are still partial. \\
        3 & Source-manifest system & Approved sources are registered with scope, URLs, update metadata, and ownership. \\
        4 & Claim-governed agent & Runtime claims, entity routing, and evidence states govern answer assembly. \\
        5 & Traceable harness & Answers, source links, validation results, and audit traces are versioned and can be replayed in validation runs. \\
        \bottomrule
    \end{tabular}
\end{table}

Applied to a new enterprise target, the pattern requires a concrete source and claim package. Public or client-approved sources must be registered, source-backed claims promoted, routing rules checked, and scenario contracts run. The public-data baseline makes these steps inspectable before private documents, production credentials, or operational logs are introduced. The framework also separates system readiness from usefulness for investment analysis: this paper evaluates preservation of source, routing, answer, trace, and hygiene contracts, while expert review and deployment evidence are left to later domain-value evaluation.

The reproducible layer consists of scenario JSON files, source and claim manifests, validation scripts, sanitized traces, and fixture or smoke-test modes. At public submission, the release package will record the public repository URL and commit hash for the submitted baseline; the current manuscript keeps this as a replaceable repository placeholder \citep{ahnkim2026advisorrepo}. Some materials cannot be redistributed directly: API keys, copyrighted PDFs, and future client-operation logs must remain private. For those materials, the reproducibility strategy provides public URLs, source manifests, checksums when available, evidence records, and claim manifests as redistributable metadata. Until the release repository is public, code and sanitized artifacts are available from the corresponding author upon request.
